\section{Scaling Laws suaves y capacidades de tarea no suaves}

La sección anterior fijó los límites de la evidencia; esta sección establece la comparación más básica. La held-out loss de los modelos de lenguaje suele seguir una power law suave en función del cómputo, los datos y los parámetros; en cambio, las métricas de tarea que de verdad interesan al usuario pueden mantenerse en chance level durante un intervalo muy largo antes de subir con rapidez. Ambas curvas pueden ser válidas al mismo tiempo, porque difieren en el objeto que miden, en la escala y en la forma de aplicar el umbral.

\subsection{Primero el mapa del artículo, después cada curva}

Las diapositivas muestran primero la portada del artículo y la lista de tareas para recordar al lector que la ``emergence'' no se generaliza a partir de un caso vistoso, sino que es un catálogo de fenómenos reunido a lo largo de varias familias de modelos, few-shot tasks y prompting techniques. El catálogo por sí solo no demuestra un mecanismo unificado, pero sí plantea preguntas de comparación reutilizables.

\lecturefigure{slide-03.jpg}{Emergent Abilities of Large Language Models: el artículo y el catálogo de capacidades}{Diapositivas oficiales, página 3; Wei et al., TMLR 2022}

Al leer esta figura, conviene distinguir primero entre «catalog» y «theory». El artículo reúne los fenómenos de umbral en tareas de arithmetic, symbolic manipulation, language understanding y otras, así como ciertas técnicas de prompting que solo funcionan en modelos más grandes; no ofrece una teoría cerrada capaz de predecir todos los umbrales directamente a partir de los parámetros y de la composición de los datos.

\subsection{Scaling law: lo predecible es la loss promedio}

La conclusión típica de las scaling laws es que, cuando las demás condiciones están controladas, la test loss disminuye aproximadamente según una ley de potencias en función de los parámetros del modelo $N$, el número de tokens de entrenamiento $D$ o el cómputo de entrenamiento $C$. En la clase se usan los tres grupos de gráficas de Kaplan et al. para establecer una línea base «suave y extrapolable», y luego se define la emergencia como el cambio a nivel de tarea que las curvas de los modelos pequeños no pueden revelar directamente. Primero hay que entender cómo cambia el error de predicción promedio; solo así se evita interpretar el abrupt-looking score de algún benchmark posterior como una contradicción con el loss scaling.

\lecturefigure{slide-04.jpg}{Mejora predecible de la loss al crecer el modelo, los datos y el cómputo}{Diapositivas oficiales, página 4; Kaplan et al., 2020}

\begin{knowledgebox}{Lectura de la figura: qué dicen las tres curvas de scaling}
Los ejes horizontales representan, respectivamente, compute, dataset size y parameter count; el eje vertical es la held-out language-model loss. Primero, compare la tendencia rectilínea global en coordenadas log--log, no un punto aislado; segundo, observe que se trata del error de next-token prediction promediado sobre la distribución, no de la exactitud en una tarea discreta; por último, una loss baja indica que el modelo predice mejor la distribución de los datos, pero no nos dice directamente si supera la random baseline de una tarea determinada.
\end{knowledgebox}

El cómputo de entrenamiento suele expresarse con una relación aproximada:
\[
C \propto N D,
\]
donde $C$ son las floating-point operations (FLOPs) del preentrenamiento, $N$ es el número de parámetros que intervienen en los cálculos hacia adelante y hacia atrás, y $D$ es el número de tokens de entrenamiento. La constante de proporcionalidad depende de la estructura del Transformer, la longitud de secuencia, si se repiten los datos de entrenamiento y los detalles de implementación; por ello no es una contabilidad exacta, sino solo una explicación de por qué, en una familia de modelos con $D$ fijo, el número de parámetros y los FLOPs de entrenamiento pueden dar ejes horizontales similares.

\teachervoice{Al inicio de la clase, por un problema al compartir la pantalla, se repitió tres veces la misma conclusión: la mejora de la loss es predecible, y la emergence de la que trata esta sesión es precisamente el cambio en las tareas que resulta difícil de predecir si solo se observan modelos pequeños. Este pequeño incidente dejó, de hecho, muy clara la diferencia entre ambos tipos de evidencia.}

\subsection{De la intuición física a una definición medible}

La intuición científica de ``More is Different'' ayuda a entender por qué un cambio cuantitativo puede producir un cambio cualitativo, pero la investigación en modelos de lenguaje necesita una operational definition reproducible. De lo contrario, cualquier salida sorprendente podría llamarse emergencia a posteriori, y el concepto perdería su valor comparativo.

\lecturefigure{slide-05.jpg}{La emergence en la ciencia: los cambios cuantitativos dan lugar a cambios cualitativos}{Diapositivas oficiales, página 5; Anderson 1972 y Jacob Steinhardt 2022}

En la clase se usan dos ejemplos intuitivos para ilustrar que «la escala cambia los estados alcanzables»: una pequeña cantidad de uranium y una cantidad suficientemente densa de uranium se encuentran en regímenes de reacción distintos; tampoco hay una simple diferencia lineal entre las moléculas pequeñas y el DNA en su capacidad para codificar información. La analogía solo sirve para construir la intuición y no sustituye los experimentos controlados en modelos de lenguaje, porque los parámetros, los datos y los algoritmos no son una magnitud escalar única como la masa física.

\lecturefigure{slide-06.jpg}{Definición operativa de las capacidades emergentes de los grandes modelos de lenguaje y los tres ejes de scale}{Diapositivas oficiales, página 6; Wei et al., 2022}

\begin{importantbox}{Definición operativa}
Si una ability no se detecta en modelos más pequeños y sí se detecta en modelos más grandes, se le llama emergent ability. Para la few-shot classification, en la clase se concreta además «no detectada» como cercana a la random accuracy, y «detectada» como significativamente superior a la random accuracy. Esta definición depende de la tarea, el prompt, la metric, la baseline y la familia de modelos que se compara.
\end{importantbox}

\begin{warningbox}{No sustituya el eje horizontal por «inteligencia»}
Los training FLOPs, el número de parámetros y la cantidad de datos de entrenamiento son tres ejes de recursos distintos. Aunque estén muy correlacionados dentro de una misma familia de modelos, ninguno de ellos puede llamarse directamente intelligence. La calidad de los datos, la token mixture, la architecture, el optimizer y el post-training pueden modificar la capacidad efectiva de modelos con el mismo número de parámetros.
\end{warningbox}

\subsection{Resumen de la sección}

Una scaling law suave describe el error de predicción promediado sobre la distribución; una emergent-ability plot describe si una tarea concreta supera la baseline en una metric concreta. La primera puede mejorar de forma continua mientras la segunda permanece estancada durante mucho tiempo, porque la exactitud en la tarea puede manifestarse solo cuando varias subhabilidades alcanzan una confiabilidad mínima.
